\chapter{结论与展望}
\markboth{结论与展望}{}

\section{全文结论}

本文围绕便携式道岔伤损检测电磁层析成像系统的设计与实现，针对道岔表面及近表面伤损检测中存在的微弱电磁响应提取困难、传统 EMT 成像分辨率有限、经典算子学习架构对铁磁 EMT 问题适配不足以及现场检测设备便携化不足等问题，从成像算法与硬件平台两个层面开展了研究，并取得了以下主要结论。

（1）在成像算法方面，本文提出了面向道岔伤损检测任务的共享条件特征驱动深度算子网络 SCF-DeepONet。针对经典 DeepONet 的 Branch--Trunk 结构难以直接处理 EMT 多激励结构化观测与连续域输出的问题，本文设计了共享条件编码、Fourier 位置编码、FiLM 调制、CNN 特征图精修以及二维--三维分支耦合机制。网络架构消融实验表明，在完全相同的训练条件和数据集下，SCF-DeepONet 相较 Vanilla DeepONet 基线在独立测试集上的平均 IoU 由 67.79\% 提升至 81.53\%，困难样本检测下限（Worst-10 Avg IoU）由 37.09\% 提升至 51.34\%。\textbf{这是本文在铁磁性 EMT 伤损检测任务上取得的最核心、最具决定性的性能提升，验证了所提出架构对强非线性、病态性逆问题的有效适配能力。}

（2）在物理信息嵌入探索方面，本文在 SCF-DeepONet 基础上进一步引入三维磁矢势差场重构作为\textbf{训练阶段辅助任务与 PINN 渐进预实验}，并设计了包含热图预训练、联合训练、热图精修和 PDE 弱物理微调的四阶段课表策略。实验表明，三维场辅助监督虽未在同分布矩形样本上取得更高的平均 IoU，但在圆形伤损零样本测试中将 Mean IoU 由 34.05\% 提升至 40.31\%，改善了模型对未见过几何形状的分布外泛化能力；基于麦克斯韦方程相对残差的 PDE 弱约束能够在不破坏二维检测性能的前提下，对三维场分支进行小幅而稳定的物理一致性修正。需要强调的是，\textbf{上述两项探索并不以提升二维检测 IoU 为核心目标，而是为物理信息神经网络在 EMT 领域的适用性提供了由浅入深的可行性验证与经验积累}。

（3）在硬件平台方面，本文设计并实现了一套便携式 EMT 系统。系统采用“测控板 + 信号切换板”的双板式结构，以 100 kHz 为主工作点，通过模拟前端窄带相敏解调将微弱交流响应转换为同相和正交两路直流量，单帧完整测量周期约为 1.46 s。基础性能测试表明系统激励稳定、锁相解调链路工作正常、通道一致性良好、互易性误差平均值约为 1.27\%。机加工仿损试块实验进一步表明，SCF-DeepONet 在实测数据上能够给出比传统线性成像方法更集中的缺陷位置响应。上述结果验证了本文所构建的“硬件采集--算法反演”端到端链路在道岔伤损检测中的工程可行性。

\section{后续研究展望}

（1）在算法层面，后续可在 SCF-DeepONet 架构基础上进一步探索多频数据融合策略，利用不同频率下趋肤深度和伤损敏感性的差异提升检测精度；同时，可考虑将 PDE 弱约束与 A 场辅助监督进行更紧密的耦合设计，而非当前分阶段独立实施的方式。

（2）在硬件层面，后续可考虑发展“模拟前端预处理 + 数字后端补偿”的混合架构，在保持现有高信噪比优势的同时，提高系统频率配置和算法适配的灵活性；此外，功耗、重量和电池续航等便携性指标也需要进一步优化。

（3）在传感器层面，本文采用的线圈分时复用结构在 $N$ 个线圈条件下只能获得 $N \times (N-1)$ 个测量通道，而根据互易定理其中真正完全独立的信息仅约为其一半，观测信息量偏少限制了反演分辨率。后续可设计异构线圈阵列、多层阵列或多频/编码激励方案，为深度学习模型提供更丰富的输入信息。另外，针对道岔伤损检测进一步需要验证柔性阵列传感器的适用性，以实现更紧贴复杂曲面结构的测量。

（4）在实验验证层面，后续应面向更接近真实工况的道岔部件和复杂伤损样本，开展更大规模实测数据采集，并进一步研究提离变化、环境干扰和结构复杂度对检测结果的影响，为 EMT 系统在真实道岔场景中的现场部署奠定基础。